

% =============================================================================================
% Introduction 

Ce court chapitre présente les bases de la logique indispensables dans le raisonnement 
mathématique, notamment en RCHI (voir plus haut). 

% =============================================================================================
% Logique Propositionnelle 

\section{Logique Propositionnelle}

\subsection{Langage Propositionnel}

Le langage propositionnel est la base du raisonnement logique. Il définit les opérateurs 
ainsi que l'ensemble des notations pour la résolution de problèmes. Il est 
défini par un \emph{alphabet}.

\begin{definition}[Alphabet et Mot]
    Un alphabet est un ensemble contenant : 
    \begin{itemize} 
        \item Un ensemble $ \mathcal{P} = (x_1, \dots, x_n, \dots)$ de variables, 
        \item Des \emph{connecteurs logiques} $ \land, \lor, \lnot, \Rightarrow, \iff$, 
        \item Et les parenthèses $ ($ et $)$. 
    \end{itemize}
    On appelle \emph{mot} de $ \mathcal{A}$ la juxtaposition d'un ensemble fini d'éléments 
    de $ \mathcal{A}$. L'ensemble des mots possibles sur $ \mathcal{A}$ est noté $ \mathcal{A}^*$. 
\end{definition}

La plupart des mots de $ \mathcal{A}$ n'ont en réalité aucun sens. Pour être capable de raisonner 
dessus, il nous faut fixer un ensemble de règles de syntaxe. 

\begin{definition}[Formule Propositionnelle]
    L'ensemble des \emph{formules propositionnelles} sur $ \mathcal{P}$, noté $ \mathcal{F}$ est le 
    plus petit sous-ensemble de mots de $ \mathcal{A}^*$ tel que : 
    \begin{enumerate}[label=\roman*)]
        \item Pour toute $ x_i \in \mathcal{ P}, x_i \in \mathcal{F}$. 
        \item Si $ P \in \mathcal{F}$, alors $ \lnot P \in \mathcal{F}$. 
        \item Si $ P, Q \in \mathcal{F}$, alors $ (P \land Q) \in \mathcal{F}$. 
        \item Si $ P, Q \in \mathcal{F}$, alors $ (P \lor Q) \in \mathcal{F}$. 
        \item Si $ P, Q \in \mathcal{F}$, alors $ P \Rightarrow Q \in \mathcal{F}$. 
    \end{enumerate}
\end{definition}





